
\subsubsection{FFWindow}

The class FFWindow is a general class that can execute tcl/tk scripts. It
defines the following methods:

\begin{itemize}

\item{\bf FFWindow()}

Creates a new instance of {\bf FFWindow}.

\item{\bf int mapWindow()}

Maps the window on the screen.

\item{\bf int unmapWindow()}

Unmaps the window from the screen.

\item{\bf int withdraw()}

Withdraws the window from the screen.

\item{\bf int iconify()}

Iconifies the Window.

\item{\bf int deiconify()}

Deiconifies the Window.

\item{\bf virtual void update()}

Does a graphic update on the window.

\item{\bf   void setDisplay(const String\& disp)}
  
  Sets the display where the window should appear to {\bf disp}.
  This method must be called before {\bf initFromFile(...)} to
  have any effect.

  void structureProc(XEvent);  
  
\item{\bf void initFromFile(const String\& file, int argc=0, char
    **argv=NULL)}
  
  Executes the tcl files {\bf file}. {\bf argc} and {\bf argv}
  are used as if they were given to the wish shell.

\item{\bf virtual int setTitle(const String\& title)}

Sets the title of the window to {\bf title}.


\end{itemize}

\subsubsection{FFGraphWindow}
\label{ffgraphwindow}

The class {\bf FFGraphWindow} is a subclass of {\bf FFGraphLabel}
and is the superclass of all window labels. Window labels are
used to display graphs. Usually you want to display a graph in a
window on a screen. But you can also subclass {\bf FFGraphWindow}
to print the graph on a printer or export it as an image or
whatever. {\bf FFGraphWindow} is a virtual class, so no instance
of this class can be created. {\bf FFGraphWindow} defines the
following methods:

\begin{itemize}

\item{\bf FFGraphWindow()}

Adds ":" followed by "FFGraphWindow" to the type of the label.

\item{\bf void scaleXYZ(FFCoord dx, FFCoord dy, FFCoord dz) =
    NULL}
  
  Should scale the coordinate system by {\bf dx} in x direction,
  by {\bf dy} in y direction and by {\bf dz} in z direction.

\item{\bf void rotateX(FFCoord dx) = NULL}
  
  Should rotate the coordinate system by {\bf dx} degrees around
  the x axis.

\item{\bf void rotateY(FFCoord dy) = NULL}
  
  Should rotate the coordinate system by {\bf dy} degrees around
  the y axis.

\item{\bf void rotateZ(FFCoord dz) = NULL}
  
  Should rotate the coordinate system by {\bf dz} degrees around
  the z axis.

\item{\bf void translate(FFCoord dx, FFCoord dy, FFCoord dz) = NULL}
  
  Should translate the coordinate system by {\bf dx} in x
  direction, by {\bf dy} in y direction and by {\bf dz} in z
  direction.
        
\item{\bf FFPoint proj(const FFPoint\& A) = NULL}
  
  Should convert the graph coordinate {\bf A} to the output
  device coordinate system. The third coordinate of the result
  should be usable by hidden-line algorithms.

\item{\bf FFPoint unProj(const FFPoint\& A) = NULL}
  
  Should convert the output device coordinate {\bf A} back to the
  graph coordinate system.

\item{\bf FFPoint adjust(FFGraphicNode *A, FFPoint B) = NULL}
  
  Should return where an edge to {\bf A} with {\bf B} being the
  last point of the edge before {\bf A} should end. E.g.  if you
  display nodes as circles you might want the edge to stop at the
  outskirts instead of going to the center of the circle.

\item{\bf void setLabelOffset(int dx, int dy) = NULL}
  
  This method sets the relative position of the name of a node to
  the position of the node.
            
\item{\bf void adjustWindowSize() = NULL}
  
  Should resize the window to display the whole graph.

\item{\bf String addNode(FFGraphicNode *n) = NULL}
  
  Should add a graphic representation of node {\bf n}. The return
  value must be a unique String by which the graphic
  representation can be addressed.

\item{\bf String addEdge(FFGraphicEdge *e) = NULL}
  
  Should add a graphic representation of edge {\bf e}. The return
  value must be a unique String by which the graphic
  representation can be addressed.

\item{\bf String addLabel(FFGraphicNode *n) = NULL}
  
  Should add a graphic representation for the name of node {\bf
    n}. The return value must be a unique String by which the
  graphic representation can be addressed.

\item{\bf String addLabel(FFGraphicEdge *e) = NULL}
  
  Should add a graphic representation for the name of edge {\bf
    e}. The return value must be a unique String by which the
  graphic representation can be addressed.
  
\item{\bf void changeColor(FFGraphicNode *n, String\& color) =
    NULL}
  
  Should change the color of node {\bf n} to {\bf color}.
  
\item{\bf String changeToIcon(FFGraphicNode *n, String\& icon) =
    NULL}
  
  Should change the representation of node {\bf n} to be
  displayed as the icon that is stored in file {\bf icon}.

\item{\bf String changeToCircle(FFGraphicNode *n) = NULL}
  
  Should change the representation of node {\bf n} to be
  displayed as a graphical item like circles, boxes etc.

\item{\bf String changeToWindow(FFGraphicNode *n) = NULL}
  
  This method is used to display nested graphs. The graphic
  representation of {\bf n} should change to something that can
  display a graph itself.

\item{\bf void changeMask(FFGraphicNode *n, String\& mask) = NULL}
  
  If node {\bf n} is displayed as an icon, the icon should be
  masked using the bitmap that is stored in file {\bf mask}.
  
\item{\bf void changeColor(FFGraphicEdge *e, String\& color) =
    NULL}
  
  The color of edge {\bf e} should change to {\bf color}.
  
\item{\bf void changeLabel(FFGraphicNode *n, const String\& name)
    = NULL}
  
  The name of node {\bf n} should be changed to {\bf name} in the
  window.
  
\item{\bf void changeLabel(FFGraphicEdge *e, const String\& name)
    = NULL}
  
  The name of edge {\bf e} should be changed to {\bf name} in the
  window.

\item{\bf FFPoint getRange(const FFGraphicNode *n) = NULL}
  
  Returns a point containing (width,height,depth) of the bounding
  box of the node {\bf n}.

\item{\bf void setSpline(FFGraphicEdge *e, BOOL spline) = NULL}
  
  Sets whether edge {\bf e} should be displayed as a spline ({\bf
    spline == YES}) or not ({\bf spline == NO}).

\item{\bf void remove(FFGraphicNode *n) = NULL}
  
  Should remove the graphic representation of node {\bf n} from
  the screen.

\item{\bf void remove(FFGraphicEdge *e) = NULL}
  
  Should remove the graphic representation of edge {\bf e} from
  the screen.

\item{\bf void removeLabel(FFGraphicNode *n) = NULL}
  
  Should remove the graphic representation of the name of node
  {\bf n} from the screen.

\item{\bf void removeLabel(FFGraphicEdge *e) = NULL}
  
  Should remove the graphic representation of the name of edge
  {\bf e} from the screen.

\item{\bf void updateGraph() = NULL}
  
  Should redraw the whole graph.

\item{\bf void updateGraph(FFGraphicNode *n) = NULL}
  
  Should redraw node {\bf n}.

\item{\bf void updateGraph(FFGraphicEdge *e) = NULL}
  
  Should redraw edge {\bf e}.

\item{\bf void setCurrent(FFGraphicNode *n) = NULL}
  
  Should unselect all items and mark node {\bf n} as selected.

\item{\bf void setCurrent(FFGraphicEdge *e) = NULL}
  
  Should unselect all items and mark edge {\bf e} as selected.

\item{\bf void addToCurrent(FFGraphicNode *n) = NULL}
  
  Should unselect any selected edges and mark node {\bf n} as
  selected.

\item{\bf void selectIfNot(FFGraphicNode *n) = NULL}
  
  Should unselect any selected edges. If node {\bf n} is selected
  unselect it, if it is not selected select it.

\item{\bf void unsetCurrent() = NULL}
  
  Should unselect all selected items.
  
\item{\bf AVLSet\verb|<|FFGraphicNode *\verb|>|\& selectedNodes()
    = NULL}
  
  Should return an AVLSet containing all selected nodes.

\item{\bf BOOL isCurrent(FFGraphicNode *n) = NULL}
  
  Should return whether node {\bf n} is selected or not.
  
\item{\bf void unset(FFGraphicNode *n) = NULL}
  
  Should unselect node {\bf n}.

\item{\bf void unset(FFGraphicEdge *e) = NULL}
  
  Should unselect edge {\bf e}.

\item{\bf void updateInspector() = NULL}
  
  Should update the inspector, if the window has one.

\item{\bf void sizeTo(int x, int y) = NULL}

Should change the size of the window to ({\bf x}, {\bf y}).
  
\item{\bf void setUpdate(BOOL f) = NULL}

Sets whether window updates shall be made.

\item{\bf BOOL getUpdate() = NULL}

Returns whether window updates are made.

\item{\bf void update() = NULL}

Forces a window update.

\item{\bf int setTitle(const String\& s) = NULL}

Set the title of the window to {\bf s}.

\item{\bf void setBusy() = NULL}
  
  Should in some way indicate that the window is busy, e.g. by
  changing the mouse cursor.

\item{\bf void unsetBusy() = NULL}
        
  Should somehow indicate that the window is not busy, e.g. by
  changing the mouse cursor.

\item{\bf int mapWindow() = NULL}

Should map the window to the display.

\item{\bf int unmapWindow() = NULL}

Should unmap the window from the display.

\item{\bf int withdraw() = NULL}

Should withdraw the window form the display.

\item{\bf int iconify() = NULL}

Should iconify the window.

\item{\bf int deiconify() = NULL}

Should deiconify the window.


\end{itemize}

\subsubsection{FFGraphicNode}
\label{ffgraphicnode}

As nodes do not contain graphic information, when {\bf
  FFGraphClass::createGraphic(...)} is invoked the graph labels
itself with the node label {\bf FFGraphicNode}. This label
contains all the graphic information that is needed to display a
node. {\bf FFGraphicNode} provides the following methods:

\begin{itemize}

\item{\bf FFGraphicNode()}
  
  Creates a new instance of {\bf FFGraphicNode}. You should never
  have to call this method, as the graph creates this label
  itself if it is needed.

\item{\bf virtual void init(FFNode n)}
  
  This method is called internally by the graph to initialize
  the label after it has been bound to a node. The label checks
  if the graph has a window label and displays itself in the
  window label if it finds one using the method {\bf
    FFGraphWindow::addNode(...)}.

\item{\bf virtual void event(labelEvent type, labelEventStruct *)}
  
  Whenever there is a structural change to the graph this method
  is invoked. This method responds to the following events:

  \begin{itemize}

  \item{\bf NAME\_CHANGE}
    
    The name on the screen is changed according to the new name
    of the node

  \item{\bf ADD\_GRAPHLABEL}
    
    Checks if the new graph label is a window label and if yes,
    it displays itself in the new window. As only one window
    label may be assigned to a graph, it is assured that the node
    is not already displayed in a window.

  \item{\bf DEL\_GRAPHLABEL}
    
    If the graph label that is going to be deleted is the window
    label the node is displayed in, the node is removed from the
    window.

  \end{itemize}

\item{\bf inline FFCoord getx() const}

  Returns the x coordinate of the node.

\item{\bf inline FFCoord gety() const}

  Returns the y coordinate of the node.

\item{\bf inline FFCoord getz() const}

  Returns the z coordinate of the node.

\item{\bf inline const FFPoint\& getCoords() const}
  
  Returns a reference to the coordinates of the node. To change
  the coordinates use one of the {\bf moveTo(...)} or {\bf
    rmoveTo} methods.

\item{\bf inline FFCoord getWidth()}

  Returns the width of the label.

\item{\bf inline FFCoord getHeight()}

  Returns the height of the label.

\item{\bf inline FFCoord getDepth()}

  Returns the depth of the label.

\item{\bf inline FFPoint\& getRange()}

  Returns the range of the label.

\item{\bf void setWidth(FFCoord w)}

  Sets the width of the node to {\bf w}.

\item{\bf void setHeight(FFCoord h)}

  Sets the width of the node to {\bf h}.

\item{\bf void setDepth(FFCoord d)}

  Sets the width of the node to {\bf d}.

\item{\bf inline void setRange(const FFPoint\& p)}

  Sets the range of the node to {\bf p}.
 
\item{\bf void changeToCircle()}
  
  Tells the node to display itself as a graphical item, like a
  circle or a box etc. . This method calls the method {\bf
    FFGraphWindow::changeToCircle(...)} to change the node. Which
  graphical item is used depends on the implementation of the
  window label.
        
\item{\bf void changeToWindow()}
  
  Tells the node to display itself as a window that may contain a
  graph.  This method calls the method {\bf
    FFGraphWindow::changeToWindow(...)}. The actual
  implementation of how to display the window is left to the
  window label.

\item{\bf const String\& getNodeWindow() const}
  
  Returns the id of the window the node is displayed as, or "" if
  the node is not displayed as a window.

\item{\bf void changeToIcon(const String\& iconFile)}
  
  Tells the node to display itself as the icon that is defined in
  {\bf iconFile}.  This method calls the method {\bf
    FFGraphWindow::changeToIcon(...)}. The actual implementation
  of how to display the icon is left to the window label.

\item{\bf String\& getIcon()}

  Returns the id of the icon the node is displayed as, or "" if
  the node is not displayed as an icon.

\item{\bf BOOL isScalable()}
  
  If the coordinate system of the window label is scaled, this
  method returns if the graphical item the node is displayed as
  should be affected by this. The method returns {\bf NO} for
  icons and window and {\bf YES}  for other items.

\item{\bf BOOL isCircle()}

  Return whether the node is displayed as a graphical item.

\item{\bf void setMask(const String\& maskFile)}
  
  Tells the node to be masked according to the mask defined in
  {\bf maskFile}. This method calls the method {\bf
    FFGraphWindow::setMask(...)}. 

\item{\bf String\& getMask()}

  Returns the id of the mask or "" if the node is not masked.

\item{\bf const String\& getNodeWindow() const}
  
  Returns the id of the window the node is displayed as, or "" if
  the node is not displayed as a window.

\item{\bf const String\& getId() const}

  Returns the id of the graphical representation of the node.

\item{\bf const String\& getLabelId() const}
  
  Returns the id of the graphical representation of the name of
  the node.

\item{\bf void setName(const String\& s)}
  
  Sets the graphical representation of the name of the node to
  {\bf s}.

\item{\bf void setColor(const String\& c)}
  
  Sets the color of the node to {\bf c}. {\bf c} must be a color-name your
  X11 server knows about (see your local rgb.txt), or specified in the
  format \verb|#|rrggbb.
  
\item{\bf String getColor()}

  Returns the color of the node.

\item{\bf void setWindow(FFGraphWindow *w)}

Displays the node in the window label {\bf }. You can use the special
window label {\bf NULL} to hide a node.

\item{\bf void moveTo(char *x, char *y, char *z="0")}
  
  Tells the node to move to the coordinates ({\bf x},{\bf y},
  {\bf z}).

\item{\bf void moveTo(const FFPoint\& p)}

  Tells the node to move to the coordinate {\bf p}.

\item{\bf void rmoveTo(char *x, char *y, char *z="0")}
  
  Tells the node to add ({\bf x},{\bf y},{\bf z}) to its current
  position.

\item{\bf void rmoveTo(const FFPoint\& p)}
  
  Tells the node to add {\bf p} to its current position.

\item{\bf void updateGraphic()}

  Tells the node to redisplay itself.

\item{\bf inline FFNodeClass *getNode() \\
 inline FFNodeClass *operator()() \\
 inline FFNodeClass *operator-\verb|>|()}

Returns the node the graphic node is bound to.

\end{itemize}

\subsubsection{FFGraphicEdge}
\label{ffgraphicedge}

As edges do not contain graphic information, when {\bf
  FFGraphClass::createGraphic(...)} is invoked the graph labels
itself with the edge label {\bf FFGraphicEdge}. This label
contains all the graphic information that is needed to display an
edge. {\bf FFGraphicEdge} provides the following methods:

\begin{itemize}

\item{\bf virtual void event(labelEvent type, labelEventStruct *)}
    
Whenever there is a structural change to the graph, this method is
called. {\bf type} specifies the type of event that occurred. The
method responds to the following methods:

\begin{itemize}

\item{\bf NAME\_CHANGE}

Changes the name of the edge according to the new name.

\item{\bf SOURCE\_CHANGE}

Redisplays the edge with the new source node.

\item{\bf TARGET\_CHANGE}

Redisplays the edge with the new target node.

\item{\bf REV\_EDGE}

Redisplays the edge with the new direction.

\item{\bf ADD\_GRAPHLABEL}

Checks whether the new label is a graph window and if it is, the edge
displays itself in the new window label. As there can only be one
window label bound to a graph it is assured that the edge is not
already displayed.

\item{\bf DEL\_GRAPHLABEL}

Checks whether the label that will be removed is a window label and
if it is the edge removes itself from the window.

\end{itemize}

\item{\bf FFCoord dist(const FFPoint\& p1, const FFPoint\& p2)}

Returns the distance between point {\bf p1} and {\bf p2}.

\item{\bf FFCoord dist(const FFPoint\& P, const FFPoint\& A, 
        const FFPoint\& B)}
  
Returns the distance between point {\bf P} and the line {\bf A} and
{\bf B}.

\item{\bf inline FFCoord getxAt(int p)}

Returns the x coordinate of the edge point with number {\bf p}

\item{\bf inline FFCoord getyAt(int p)}

Returns the y coordinate of the edge point with number {\bf p}

\item{\bf inline FFCoord getzAt(int p)}

Returns the z coordinate of the edge point with number {\bf p}

\item{\bf int getPointNear(const FFPoint\& p)}

Returns the edge point that is nearest to point {\bf p}.

\item{\bf BOOL addPointAt(int k, const FFPoint\& p)}

Add an edge point as number {\bf k} and position {\bf p}. The number
of all following edge points is increased by one.

\item{\bf BOOL deletePointAt(int k)}

Deletes the edge point with number {\bf k}.

\item{\bf BOOL deletePointAt(const FFPoint\& p)}

Deletes the point that is nearest to point {\bf p}.

\item{\bf int anzPoints()}

Returns the number of edge points.

\item{\bf FFPoint\& getPointAt(int k)}

Returns the point with number {\bf k}.

\item{\bf const MPlex\verb|<|FFPoint\verb|>|\& getEdgePoints()}
                        
Returns a MPlex containing all edge points.

\item{\bf String getId()}

Returns the id of the graphical representation of the edge.

\item{\bf String getLabelId()}

Returns the id of the graphical representation of the name of the
edge.

\item{\bf void revert()}

Reverts the edge in the window.

\item{\bf void setName(const String\& s)}

Sets the name of the edge in the window to {\bf s}.

\item{\bf virtual void init(FFEdgeClass *)}

Is called internally by the library to initialize the label. The label
Checks whether a window label is bound to the graph and displays
itself in the window if it find one.
        
\item{\bf void updateGraphic()}

Redisplays the edge.
  
\item{\bf void toggleSpline()}

Toggles whether the edge is displayed as a spline or not.

\item{\bf void setColor(const String\& c)}

Sets the color of the node to {\bf c}.

\item{\bf String getColor()}

Returns the color of the edge.

\item{\bf void rmoveTo(const FFPoint\& p)}

Adds {\bf p} to the positions of all edge points.

\item{\bf void rmovePoint(int k, const FFPoint\& p)}

Adds {\bf p} to the edge point with number {\bf k}.

\item{\bf void movePoint(int k, const FFPoint\& p)}

Moves the edge point with number {\bf k} to the position {\bf P}.

\item{\bf inline FFEdgeClass *getEdge()\\
 inline FFEdgeClass *operator()()\\
 inline FFEdgeClass *operator\verb|->|()}

Returns the edge the label is bound to.

\end{itemize}

\subsubsection{FFInspector}

Graph labels may contain inspectors to display their content. The
class {\bf FFInspector} defines a general interface for
inspectors. The class defines the following methods:

\begin{itemize}

\item{\bf FFInspector(FFGraphLabel *g)}

Creates a new instance of FFInspector and initializes it according to
the graph-label {\bf g}.

\item{\bf virtual void show()}

Should display the inspector on the screen.

\item{\bf virtual void hide()}

Should remove the inspector from the screen.

\item{\bf virtual void inspect() = NULL}

Should display the content of the graph-label.

\end{itemize}

\subsubsection{FFTclInspector}

{\bf FFTclInspector} is a superclass for all inspectors that use the
tcl/tk widget set to display themselves. It implements the following
methods.

\begin{itemize}

\item{\bf virtual void hide()}

Calls a tk function to remove the inspector from the screen.

\item{\bf virtual void show()}

Calls a tk function to show the inspector on the screen and top it.

\end{itemize}

\subsubsection{FFSugiInsp}

The class {\bf FFSugiInsp} is an inspector for the FFSugi
layout-label. It enables you to change the parameters of the sugyjama
layout algorithm.

\begin{figure}[htb]
\fbox{\epsfig{figure=sugiInsp.eps}}
\caption{FFSugiInsp}
\end{figure}
\clearpage

It implements the following methods:

\begin{itemize}

\item{\bf FFSugiInsp(FFGraphLabel *l)}

Creates a new instance of the class {\bf FFSugiInsp} and initializes
it with the sugiyama label {\bf l}.

\item{\bf int recall(Tcl\_Interp *i, int k, char **c)}

Is called by the tk library whenever the user changes any parameters
in the window.

\item{\bf void inspect()}

Displays the contents of the sugiyama layout label in its window.

\end{itemize}

\subsubsection{FFSpringInsp}

The class {\bf FFSpringInsp} is an inspector for the FFSpring
layout-label. It enables you to change the parameters of the
spring-embedder layout algorithm.

\begin{figure}[htb]
\mbox{\epsfig{figure=springInsp.eps}}
\caption{FFSpringInsp}
\end{figure}
\clearpage

It implements the following methods:

\begin{itemize}

\item{\bf FFSpringInsp(FFGraphLabel *l)}

Creates a new instance of the class {\bf FFSpringInsp} and initializes
it with the spring-embedder label {\bf l}.

\item{\bf int recall(Tcl\_Interp *i, int k, char **c)}

Is called by the tk library whenever the user changes any parameters
in the window.

\item{\bf void inspect()}

Displays the contents of the spring-embedder layout label in its
window.

\end{itemize}

\subsubsection{FFTclBasicWindow}

The class {\bf FFTclBasicWindow} is a subclass of {\bf FFGraphWindow}
and {\bf FFWindow}. It is the superclass for all window-labels that
use the tcl/tk widget set to display graphs. The class provides the
following methods:

\begin{itemize}

\item{\bf virtual void event(labelEvent type, labelEventStruct *)}

This method is invoked whenever there is a structural change in the
graph.  {\bf type} specifies which type of event occurred. The class
responds to the following methods:

\begin{itemize}

\item{\bf DIRECTEDNESS\_CHANGE}

The graph is redisplayed using arrowheads at edges if the graph is
directed.

\end{itemize}

\item{\bf virtual void scaleXYZ(FFCoord x, FFCoord y, FFCoord z)}

Scales the x-coordinate of the window coordinate system {\bf x} times,
the y-coordinate {\bf y} times and the z-coordinate {\bf z} times. 

\item{\bf virtual void rotateX(FFCoord x)}

Scales the x-coordinate of the window coordinate system {\bf x} times.

\item{\bf virtual void rotateY(FFCoord y)}

Scales the y-coordinate of the window coordinate system {\bf y} times.

\item{\bf virtual void rotateZ(FFCoord z)}

Scales the z-coordinate of the window coordinate system {\bf z} times.

\item{\bf virtual void translate(FFCoord x, FFCoord y, FFCoord z)}

Translates the window coordinate system by ({\bf x},{\bf y},{\bf z}).
        
\item{\bf virtual FFPoint adjust(FFGraphicNode *A, FFPoint B)}
  
{\bf B} is the last edge point of an edge before it reaches {\bf A}.

The result of the method is the position where the edge should meet
{\bf A}.

\item{\bf virtual BOOL hasOwnInspector() const}

Return {\bf NO}.

\item{\bf virtual void init(const FFGraph\& g)}

Initializes the window label with graph {\bf g}. The  layout labels
{\bf FFSugi} and {\bf FFSpring} are added to the graph.

\item{\bf virtual void adjustWindowSize()}

Adjusts the size of the canvas the graph is displayed in to be able
to display the whole graph.

\item{\bf virtual void setLabelOffset(int x, int y)}
  
Sets the relative position of the name of a node to the position of the
node.

\item{\bf virtual String addNode(FFGraphicNode *n)}

Adds a graphical representation for the node label {\bf n} and returns
a unique id.

\item{\bf virtual String addEdge(FFGraphicEdge *e)}

Adds a graphical representation for the edge label {\bf e} and returns
a unique id.

\item{\bf virtual String addLabel(FFGraphicNode *n)}

Adds a graphical representation for node {\bf n}.

\item{\bf virtual String addLabel(FFGraphicEdge *e)}

Adds a graphical representation for edge {\bf e}.

\item{\bf virtual void changeColor(FFGraphicNode *n, String\& color)}

Changes the color of node {\bf n} to {\bf color}.

\item{\bf virtual String changeToIcon(FFGraphicNode *n, 
        String\& iconFile)}

Displays node {\bf n} as the icon defined in {\bf iconFile}. Supported
graphic types are {\bf gif}, {\bf ppm}, {\bf xbm} and {\bf xpm}. {\bf
iconFile} must be located in the according sub-directory of the icon
directory you specified when installing the library. If e.g. {\bf
iconFile} is name.xxx there must be two files,
ICON\_LIBDIR/xxx/nameN.xxx and ICON\_LIBDIR/xxx/nameS.xxx where
nameN.xxx contains the icon for the node if it is not selected and
nameS.xxx contains the icon for the node if it is selected.

\item{\bf virtual String changeToCircle(FFGraphicNode *n)}

Displays node {\bf n} as a circle.

\item{\bf virtual String changeToWindow(FFGraphicNode *n)}

Displays node {\bf n} as a sub-window which may display another graph.

\item{\bf virtual void changeMask(FFGraphicNode *n, String\& maskFile)}

If node {\bf n} is displayed as an icon and the graphic type of the
icon is {\bf xpm} or {\bf xbm}, the icon is masked according to the
bitmap defined in {\bf maskFile}. {\bf maskFile} must be located in
the xbm subdirectory of the icon-directory you specified when
installing the library.

\item{\bf virtual void changeColor(FFGraphicEdge *e, String\& color)}

Changes the color of edge {\bf e} to {\bf color}.

\item{\bf virtual void changeLabel(FFGraphicNode *n, const 
        String\& name)}

Changes the graphical representation of the name of node {\bf n} to
display the string {\bf name}.

\item{\bf virtual void changeLabel(FFGraphicEdge *n, const 
        String\& name)}

Changes the graphical representation of the name of edge {\bf e} to
display the string {\bf name}.

\item{\bf virtual void remove(FFGraphicNode *n)}

Removes the node {\bf n} from the window.

\item{\bf virtual void remove(FFGraphicEdge *e)}

Removes the edge {\bf e} from the window.

\item{\bf virtual void removeLabel(FFGraphicNode *n)}

Removes the graphical representation of the name of node {\bf n} from
the window.

\item{\bf virtual void removeLabel(FFGraphicEdge *e)}

Removes the graphical representation of the name of edge {\bf e} from
the window.

\item{\bf virtual void updateGraph()}

Redisplays the whole graph.

\item{\bf virtual void updateGraph(FFGraphicNode *n)}

Redisplays node {\bf n}.

\item{\bf virtual void updateGraph(FFGraphicEdge *e)}

Redisplays edge {\bf e}.

\item{\bf virtual void setSpline(FFGraphicEdge *e, BOOL flag)}

If {\bf flag == YES} edge {\bf e} is displayed as a spline. If {\bf
flag == NO} the edge is displayed as lines.
  
\item{\bf virtual void setCurrent(FFGraphicNode *n)}

Unselects any selected items and selects node {\bf n}.

\item{\bf virtual void setCurrent(FFGraphicEdge *e)}

Unselects all selected items and selects edge {\bf e}.

\item{\bf virtual void addToCurrent(FFGraphicNode *n)}

Unselects all selected edges and selects node {\bf n}.

\item{\bf virtual void selectIfNot(FFGraphicNode *n)}

Unselects node {\bf n} if it is selected or call {\bf addToCurrent(n)}
if is not selected.

\item{\bf virtual void unsetCurrent()}

Unselects any selected items.

\item{\bf virtual BOOL isCurrent(FFGraphicNode *n)}

Returns whether node {\bf n} is selected.

\item{\bf virtual AVLSet\verb|<|FFGraphicNode *\verb|>|\& 
        selectedNodes()}

Returns a reference to the set containing all selected nodes.

\item{\bf virtual void unset(FFGraphicNode *n)}

If node {\bf n} is selected, it is unselected.

\item{\bf virtual void unset(FFGraphicEdge *e)}

If edge {\bf e} is selected, it is unselected.

\item{\bf virtual int recall(Tcl\_Interp *i, int k, char **c)}

This method is called by the tk library whenever the user does any
input in the window.

\item{\bf const Tcl\_Interp *getInterp()}

Returns the tk interpreter the class uses.

\item{\bf virtual void sizeTo(int w, int h)}

Resizes the window to width {\bf w} and height {\bf h}.

\item{\bf FFPoint getRange(const FFGraphicNode *n)}

Returns the range of the bounding box of {\bf n}.
  
\item{\bf virtual void setUpdate(BOOL f)}

Sets whether graphic updates should be made in the window.

\item{\bf virtual BOOL getUpdate()}

Returns whether graphic updates are made in the window.

\item{\bf virtual void update()}

Forces a graphic update in the window.
  
\item{\bf virtual int setTitle(const String\& s)}

Sets the title of the window to {\bf s}.

\item{\bf virtual void setBusy()}

Changes the mouse cursor to display a watch.

\item{\bf virtual void unsetBusy()}

Changes the mouse cursor to display an arrow,

\item{\bf virtual int mapWindow()}

Maps the window to the screen.

\item{\bf virtual int unmapWindow()}

Unmaps the window from the screen.

\item{\bf virtual int withdraw()}

Withdraws the window from the screen.

\item{\bf virtual int iconify()}

Iconifies the window.

\item{\bf virtual int deiconify()}

Deiconifies the window.

\item{\bf String fsBox(String message= String("Select File:"), 
               String type= String("*"), String name=String(""))}

Displays a file-selector box. {\bf message} contain the title of the
box, {\bf type} contains an extension filter and {\bf name} may
contain a directory or file that should be preselected.

\item{\bf FFTclBasicWindow()}

Creates a new instance of class {\bf FFTclBasicWindow}. You usually
do not use this class directly, but a subclass of this class.

\end{itemize}

\subsubsection{FFTclWindow}

The class {\bf FFTclWindow} is a subclass of {\bf FFTclBasicWindow}
and is able to display a graph in a window on the screen using the
tcl/tk widget library. The class contains an inspector of the class
{\bf FFTclGraphInsp} to display the content of the labels bound to the
graph. {\bf FFTclWindow} provides the following methods:

\begin{itemize}

\item{\bf FFTclWindow()}

Creates a new instance of the class {\bf FFTclWindow}. This instance
may be bound to a graph using the {\bf
FFGraphClass::createGraphic(FFGraphWIndow *)} method to display the
graph on the screen.

\item{\bf virtual void init(const FFGraph\& g)}

Initializes the window with graph {\bf g} and puts a window 
displaying {\bf g} on the screen.

\item{\bf virtual void event(labelEvent type, labelEventStruct *)}

This method is invoked from the library if there is any structural
change to the graph. {\bf type} specifies the type of change. The
method responds to the following event types:

\begin{itemize}

\item{\bf NAME\_CHANGE}

Sets the title of the window according to the new name of the graph.

\item{\bf ADD\_NODELABEL\\
        DEL\_NODELABEL\\
        ADD\_EDGELABEL\\
        DEL\_EDGELABEL} 

Updates the inspector.

\end{itemize}

\item{\bf virtual int recall(Tcl\_Interp *i, int k, char **c)}

This method is invoked from the tcl/tk library whenever the user makes
any input in the window.

\item{\bf virtual int exportSelection(int , char*, int)}

This method is invoked by the tcl/tk library. It exports the graph as
a string containing the graph in sgraph format.

\item{\bf virtual int getSelection(char* str)}

This method is invoked when the user does a copy or paste. It builds a
string containing the current graph in sgraph format.

%  virtual void setCurrent(FFGraphicNode *);
%  virtual void setCurrent(FFGraphicEdge *);
%  virtual void addToCurrent(FFGraphicNode *);

\item{\bf virtual void updateInspector()}

Updates the inspector.

\item{\bf virtual void showInspector()}

Brings the inspector on the screen.
  
\end{itemize}

\subsubsection{FFTclSubWindow}

This class can display a graph in a node which is displayed as a
window in the {\bf FFTclWindow} window label. The class provides 
the following methods:

\begin{itemize}

\item{\bf virtual FFGraphLabel *proto() const}

Returns a new instance of {\bf FFTclSubWindow}. This method is used
internally by the library.
  
\item{\bf virtual int recall(Tcl\_Interp *i, int k, char **c)}

This method is called by the tcl/tk library whenever the user triggers
an event in the sub-window.

\item{\bf virtual void init(const FFGraph\& g)}

Initializes the class with graph {\bf g}.
       
\item{\bf FFTclSubWindow(FFNode dad= NULL, Tcl\_Interp *i=NULL,
        Tk\_Window win=0, String\& w=String(""), String\&
        id=String(""))}

Creates a new instance of {\bf FFTclSubWindow}. {\bf dad} is the node
the subgraph of this window is bound to. {\bf i} and {\bf win} are the
tcl interpreter and the window where the subgraph should be displayed
in.  q{\bf w} is the name of the canvas widget where the subgraph
should be displayed in and {\bf id} is the tk item id of the sub-graph
window.

\end{itemize}

\subsubsection{FFPoint}

The class {\bf FFPoint} implements 3 dimensional vectors to store
coordinates. The class actually contains 4 dimensions where the fourth
dimension is necessary to be able to do translations with linear
operators. The class defines the usual operations expected for
numeric data-types. It also defines {\bf FF\_EPSILON} to 0.000000001. If the
distance between two {\bf FFPoints} is smaller than {\bf FF\_EPSILON} they are
considered to be equal. You can access a single dimension by reading
the member variable v\verb|[]|. Methods without description do what
you would expect them to do.

\begin{itemize}

\item{\bf FFPoint()}

Create a new instance and initializes it with (0,0,0).

\item{\bf FFPoint(const FFPoint\& p)}
\item{\bf FFPoint(FFCoord nx, FFCoord ny, FFCoord nz=0.0)}

\item{\bf BOOL operator == (const FFPoint\& p) const}

Returns {\bf YES} if the distance between the points is smaller than
{\bf FF\_EPSILON}.

\item{\bf FFPoint operator-() const}

\item{\bf FFPoint operator + (const FFPoint\& p) const}
\item{\bf FFCoord operator * (const FFPoint\& p) const}
\item{\bf FFPoint operator * (FFCoord i) const}
\item{\bf FFPoint operator / (FFCoord i) const}

\item{\bf void operator -= (const FFPoint\& p)}
\item{\bf void operator += (const FFPoint\& p)}
\item{\bf void operator *= (FFCoord i)}
\item{\bf void operator /= (FFCoord i)}

\item{\bf FFPoint\& operator = (const FFPoint\& v)}
\item{\bf char * operator()() const}

Returns an ASCII representation of the point.

\item{\bf FFCoord abs() const}

Returns the absolute value of the point.

\item{\bf FFCoord abs2D() const}

Returns the absolute value of the first two dimensions of the point.

\end{itemize}

\subsubsection{FFMatrix}

The class {\bf FFMatrix} implements a 3x3 Matrix. It is actually a 4x4
Matrix, but the fourth dimension is used to be able to do translations
with linear operators. The class implements the methods you expect
from a numeric data-type. You can access a row by reading
the member variable row\verb|[]|. Methods without description do what
you would expect them to do.

\begin{itemize}

\item{\bf FFMatrix()}

Creates a 0-Matrix.

\item{\bf FFMatrix(const FFPoint\&, const FFPoint\&, 
        const FFPoint\&, const FFPoint\&)}
    
\item{\bf BOOL operator == (const FFMatrix\&) const}

\item{\bf FFMatrix operator-() const}
\item{\bf friend FFMatrix operator - (const FFMatrix\&, const 
        FFMatrix\&)}

\item{\bf FFMatrix operator + (const FFMatrix\&) const}
\item{\bf FFMatrix operator* (const FFMatrix\&) const}
\item{\bf FFMatrix operator* (FFCoord) const}
\item{\bf FFPoint operator* (const FFPoint\&) const}
\item{\bf FFMatrix operator / (FFCoord) const}

\item{\bf void operator -= (const FFMatrix\&)}
\item{\bf void operator += (const FFMatrix\&)}
\item{\bf void operator *= (FFCoord)}
\item{\bf void operator *= (const FFMatrix\&)}
\item{\bf void operator /= (FFCoord)}

\item{\bf char * operator()() const}

Returns an ASCII representation of the Matrix.

\end{itemize}
